% C057 | 三种证明：最快 / 门槛最低 / 常规韦达（相互独立）
\begin{proofbox}
  \textbf{\textcolor{CProof}{先证一个零门槛小工具：焦半径公式}}

  只证右焦点 $F(c,0)$（左焦点把图左右镜像即可）。双曲线上任取 $P(x,y)$，由
  $y^2=\dfrac{b^2}{a^2}x^2-b^2$、$c^2=a^2+b^2$ 得
  \[
    \begin{aligned}
      |PF|^2&=(x-c)^2+y^2\\
      &=\frac{c^2}{a^2}x^2-2cx+a^2
       =\left(\frac ca x-a\right)^2.
    \end{aligned}
  \]
  设 $e=c/a>1$、$p=b^2/a=ec-a$，于是
  \[
    \boxed{|PF|=|ex-a|}
    \quad\Longrightarrow\quad
    \begin{cases}
      |PF|=ex-a,&P\text{ 在右支},\\
      |PF|=a-ex,&P\text{ 在左支}.
    \end{cases}
  \]
  以下三法都只使用这一个工具（常规法再使用韦达）。

  \medskip
  \textbf{\textcolor{CProof}{证明一（最快：方向余弦，一眼消去角度）}}

  \textbf{切入点：} 两个点都在\textbf{经过 $F$ 的同一条直线}上；以 $F$ 为起点，方向只影响 $\cos\theta$，目标是把它消掉。

  从 $F(c,0)$ 沿某射线走 $r=|FP|$，方向与 $x$ 正向夹角 $\theta$，则 $x=c+r\cos\theta$。记 $u=\cos\theta$。
  \begin{description}[style=nextline,leftmargin=2.8em,labelsep=0.5em,itemsep=0.4em]
    \item[\textbf{一步：写出射线上的倒数。}]
    若 $P$ 在右支，$r=e(c+ru)-a=p+eru$；若在左支，$r=a-e(c+ru)=-p-eru$。因此
    \[
      \boxed{\text{右支： }\frac1r=\frac{1-eu}{p}},\qquad
      \boxed{\text{左支： }\frac1r=\frac{-1-eu}{p}}.
    \]
    \item[\textbf{同支：相反方向，直接加。}]
    两点必在右支，且 $F$ 处于它们中间，两射线的余弦分别为 $u,-u$：
    \[
      \frac1{|AF|}+\frac1{|BF|}
      =\frac{1-eu}{p}+\frac{1+eu}{p}
      =\boxed{\frac2p=\frac{2a}{b^2}}.
    \]
    \item[\textbf{异支：同一方向，直接减。}]
    设 $A$ 为右支、$B$ 为左支，两点都在 $F$ 指向左侧的同一射线上，共用一个 $u$：
    \[
      \frac1{|AF|}-\frac1{|BF|}
      =\frac{1-eu}{p}-\frac{-1-eu}{p}
      =\boxed{\frac2p=\frac{2a}{b^2}}.
    \]
  \end{description}
  \textbf{优势：} 全程不设直线斜率、不解交点、不用韦达；还自动证明异支时右支上的点离右焦点更近。

  \medskip
  \textbf{\textcolor{CProof}{证明二（最低门槛：只靠共线比例，无三角函数）}}

  \textbf{为什么设距离？} 题目要的是 $|AF|,|BF|$ 的倒数，所以\textbf{直接设} $r_A=|AF|,\ r_B=|BF|$；不设 $A(x_1,y_1),B(x_2,y_2)$ 四个未知量。焦半径公式立刻给出各点横坐标：
  \[
    \begin{array}{lll}
      \text{右支：}& r=ex-a,&\displaystyle x-c=\frac{r-p}{e};\\[4pt]
      \text{左支：}& r=a-ex,&\displaystyle x-c=\frac{-r-p}{e}.
    \end{array}\tag{*}
  \]
  这里仅用到 $ec-a=p$。又因为 $A,F,B$ 共线，\textbf{两条直角三角形的横边与斜边成比例}；同向取等号，反向取相反数。
  \begin{description}[style=nextline,leftmargin=2.8em,labelsep=0.5em,itemsep=0.45em]
    \item[\textbf{同支（$A-F-B$）：}]
    $A,B$ 都在右支，分别把 $(*)$ 代入“\textbf{方向相反}”的关系：
    \[
      \frac{x_A-c}{r_A}=-\frac{x_B-c}{r_B}
      \quad\Longrightarrow\quad
      \frac{r_A-p}{e r_A}=-\frac{r_B-p}{e r_B}.
    \]
    两边乘以 $e$ 并移项：
    \[
      1-\frac p{r_A}=-1+\frac p{r_B}
      \quad\Longrightarrow\quad
      \boxed{\frac1{r_A}+\frac1{r_B}=\frac2p}.
    \]
    \item[\textbf{异支（$F-A-B$）：}]
    $A$ 在右支、$B$ 在左支，使用“\textbf{方向相同}”的关系：
    \[
      \frac{x_A-c}{r_A}=\frac{x_B-c}{r_B}
      \quad\Longrightarrow\quad
      \frac{r_A-p}{e r_A}=\frac{-r_B-p}{e r_B}.
    \]
    同样只需移项：
    \[
      1-\frac p{r_A}=-1-\frac p{r_B}
      \quad\Longrightarrow\quad
      \boxed{\frac1{r_A}-\frac1{r_B}=\frac2p}.
    \]
  \end{description}
  \textbf{优势：} 只需要\textbf{设两段长度、焦半径公式、共线比例}三件事；不用记三角函数和二次方程。核心图像是\textbf{同支 $A-F-B$、异支 $F-A-B$}。

  \medskip
  \textbf{\textcolor{CProof}{证明三（常规：直线联立 + 韦达，完整代数核验）}}

  设 $F(c,0)$。先取斜率存在的直线 $y=k(x-c)$，\textbf{且 $k\neq\pm b/a$}（否则仅有一个有限交点）。代入双曲线整理：
  \[
    (b^2-a^2k^2)x^2+2a^2ck^2x-a^2(c^2k^2+b^2)=0.
  \]
  两根 $x_1,x_2$ 对应两个有限交点，韦达给出
  \[
    x_1+x_2=\frac{2a^2ck^2}{a^2k^2-b^2},\qquad
    x_1x_2=\frac{a^2(c^2k^2+b^2)}{a^2k^2-b^2}.
  \]
  为避免一边写 $|ex_i-a|$ 一边反复拆绝对值，先计算\textbf{带符号的} $u_i=ex_i-a$：
  \[
    \begin{aligned}
      N:=u_1+u_2&=e(x_1+x_2)-2a
      =\frac{2ab^2(1+k^2)}{a^2k^2-b^2},\\
      M:=u_1u_2&=e^2x_1x_2-ae(x_1+x_2)+a^2
      =\frac{b^4(1+k^2)}{a^2k^2-b^2}.
    \end{aligned}
  \]
  因而（分母不为零）
  \[
    \boxed{\frac1{u_1}+\frac1{u_2}=\frac{u_1+u_2}{u_1u_2}
    =\frac NM=\frac{2a}{b^2}}.\tag{**}
  \]
  \begin{description}[style=nextline,leftmargin=2.8em,labelsep=0.5em,itemsep=0.4em]
    \item[\textbf{同支：}] 两点均在右支，$u_1,u_2>0$，于是 $r_A=u_1,r_B=u_2$。由 $(**)$ 得 $1/r_A+1/r_B=2a/b^2$。
    \item[\textbf{异支：}] 若 $A$ 在右支、$B$ 在左支，则 $u_A>0,u_B<0$；因此 $r_A=u_A,r_B=-u_B$，由 $(**)$ 得 $1/r_A-1/r_B=2a/b^2$，交换标签后写绝对值即可。
    \item[\textbf{斜率不存在：}] $x=c$ 时，$y=\pm b^2/a$，两交点均在右支；$r_A=r_B=b^2/a=p$，故 $1/r_A+1/r_B=2/p$。
    \item[\textbf{渐近线方向：}] $k=\pm b/a$ 时代入后的\textbf{二次项消失}，方程退化为一次方程，只得到一个有限交点；这两种情况不在定理的“两交点”前提内。
  \end{description}
  左焦点 $F(-c,0)$ 的结果由关于 $y$ 轴的对称性立即得到。

  \medskip
  \textbf{\textcolor{CProof}{三法选择}}

  \textbf{考场秒证：} 证明一；\quad
  \textbf{第一次理解、最低依赖：} 证明二；\quad
  \textbf{标准解析几何书面推导：} 证明三。
\end{proofbox}
